\section{理论与推导}

\subsection{系统 MTF 与 SR 可行性分析}

本节从系统的调制传递函数（\eng{MTF}）和信噪比（\eng{SNR}）出发，建立 \(2\times\) 与 \(4\times\) 超分辨率的理论可行性边界，
为主文~\S3.1 和~\S1 的 claim 边界提供定量支撑。

\subsubsection{系统 MTF 模型}

本系统的光学 \eng{PSF} 建模为各向同性 Gaussian，标准差 \(\sigma\) 以 \eng{LR} 像素为单位（1 \eng{LR} px \(=\) 10\,\(\mu\)m）。对应的 \eng{MTF} 为
\[
  \mathrm{MTF}_{\mathrm{PSF}}(f) = \exp\!\bigl(-2\pi^2 \sigma^2 f^2\bigr),
\]
其中 \(f\) 为空间频率（cyc/\eng{LR} px）。探测器孔径的 \eng{MTF} 为
\[
  \mathrm{MTF}_{\mathrm{det}}(f) = \bigl|\mathrm{sinc}(W f)\bigr|, \quad W = 10\;\mu\mathrm{m}.
\]
系统总 \eng{MTF} 为二者之积。\eng{EP15 M3} 的 \eng{FRC} 形状拟合使用了 \(\mathrm{MTF}^2\) 形式，即
\[
  \mathrm{MTF}(f)^2 = \bigl[\exp\!\bigl(-2\pi^2\sigma_{\mu\mathrm{m}}^2 f^2\bigr) \cdot \bigl|\mathrm{sinc}(10f)\bigr|\bigr]^2, \quad \sigma_{\mu\mathrm{m}} = \sigma_{\mathrm{LR}} \times 10\;\mu\mathrm{m}.
\]
需要注意的是，主文~\S3.1 引用的 \eng{MTF} 数值表仅含 Gaussian \eng{PSF} 项，未乘 detector sinc；\eng{EP15 M3} 拟合使用的则是含 sinc 的系统 \(\mathrm{MTF}^2\)。下文分别标明各自口径。

当前系统的已校准空间分辨率为 20\,\(\mu\)m，探测器采样间距（pitch）为 10\,\(\mu\)m/px，对应 \eng{FOV} 6.4\(\times\)4.8\,mm。
在 \(1\times/2\times/4\times\) 输出网格上，Nyquist 频率分别为 0.5、1.0、2.0\,cyc/detector px。下表给出 Gaussian-only 口径下各网格 Nyquist 处的 \eng{MTF} 值。

\begin{table}[ht]
  \centering
  \caption{各输出网格 Nyquist 频率处的 Gaussian-only \eng{MTF} 值}
  \label{tab:mtf-nyquist}
  \small
  \begin{tabular}{lcccc}
    \toprule
    网格 & \(f\) (cyc/px) & \(\sigma=0.20\) & \(\sigma=0.35\) & \(\sigma=0.50\) \\
    \midrule
    \(1\times\) Nyquist & 0.5 & 0.821 & 0.546 & 0.291 \\
    \(2\times\) Nyquist & 1.0 & 0.454 & 0.089 & 0.007 \\
    \(4\times\) Nyquist & 2.0 & 0.042 & \(6.3\times10^{-5}\) & \({\sim}3\times10^{-9}\) \\
    \bottomrule
  \end{tabular}
\end{table}

\subsubsection{有效 SNR 判据}

为判断某一空间频率处的信息是否可恢复，定义有效 \eng{SNR} 为
\[
  \mathrm{SNR}_{\mathrm{eff}}(f) = \frac{\Delta T \cdot \mathrm{MTF}(f,\sigma)}{\sigma_n},
\]
其中 \(\Delta T\) 为热对比度，\(\sigma_n = 0.0724\;^\circ\mathrm{C}\) 为探测器噪声底（由平滑区域相邻坐标 \eng{MAE} 测定）。
实测对比度覆盖从噪声底到外轮廓中位 2.49\,\(^\circ\)C 的宽范围：名义边缘 0.70\,\(^\circ\)C（输入 \eng{SNR} \(=\) 9.7）、内轮廓中位 1.94\,\(^\circ\)C（\eng{SNR} \(=\) 26.8）、外轮廓中位 2.49\,\(^\circ\)C（\eng{SNR} \(=\) 34.4）。

将上述对比度代入各网格 Nyquist 处 \eng{MTF}，得到有效 \eng{SNR} 表。

\begin{table}[ht]
  \centering
  \caption{各对比度在 \(2\times\)/\(4\times\) Nyquist 处的有效 \eng{SNR}}
  \label{tab:snr-eff}
  \small
  \begin{tabular}{lcc}
    \toprule
    对比度 & \(2\times\)（\(\sigma\)=0.20 / 0.35 / 0.50） & \(4\times\)（\(\sigma\)=0.20 / 0.35 / 0.50） \\
    \midrule
    名义 0.7\,\(^\circ\)C    & 4.39 / 0.86 / 0.07   & 0.41 / 0.001 / \(\sim\)0 \\
    内轮廓 1.94\,\(^\circ\)C & 12.16 / 2.39 / 0.19  & 1.14 / 0.002 / \(\sim\)0 \\
    外轮廓 2.49\,\(^\circ\)C & 15.62 / 3.06 / 0.25  & 1.46 / 0.002 / \(\sim\)0 \\
    \bottomrule
  \end{tabular}
\end{table}

风险带定义为：\(\mathrm{SNR}_{\mathrm{eff}} \ge 5\) 为 observable，\(\ge 3\) 为 borderline，\(\ge 1\) 为 weak，其余为 noise-dominated。

\subsubsection{可行性判定}

\textbf{\(2\times\) 在条件下可行。} 在乐观 \(\sigma\) 端（\(\sigma \le 0.35\)），内/外轮廓的 \(\mathrm{SNR}_{\mathrm{eff}}\) 处于 observable 到 borderline 区间，
存在可用的频率余量。但当 \(\sigma = 0.5\) 时，\(2\times\) Nyquist 处 \eng{MTF} 仅剩 0.007，高频信号强烈衰减至噪声以下。
因此本 \eng{POC} 将目标定位在 contour-level 增强，而非计量级温度恢复。

\textbf{\(4\times\) 出界。} 除最乐观的 \(\sigma = 0.2\) 在外轮廓对比度下勉强达到 weak（1.46），其余条件下 \(\mathrm{SNR}_{\mathrm{eff}}\) 全线 noise-dominated。
\(4\times\) 网格仅作 contour oversampling 和可视化用途。这一理论预测与后续 \(4\times\) 网络实验的负结果一致（见 D.4.2）——理论界先于实验设定了风险预期。

需要强调的是，\(\mathrm{SNR}_{\mathrm{eff}}\) 过门是必要条件而非充分证明：该判据忽略了对齐误差、热漂移与模型失配等因素。

作为辅助证据，Cramér-Rao 下界（\eng{CRB}）给出了边缘定位精度的理论极限。例如在 \(\Delta T = 0.7\;^\circ\mathrm{C}\)、\(\sigma_{\mathrm{PSF}} = 0.5\;\mathrm{px}\) 条件下，
单帧 \eng{CRB} 为 0.1415\,px，16 帧平均后降至 0.0344\,px。\eng{EP04} 实测 A-class split-half 约 0.027\,px 与 16 帧 \eng{CRB} 同量级，二者作为一致性检验互相支撑。

\subsection{观测算子与零空间}

\subsubsection{前向模型与记号}

主文~\S3.1 定义的离散观测模型为
\[
  y_k = D \cdot B \cdot H \cdot S_{t_k} \cdot x + n_k,
\]
其中 \(x\) 为 \(2\times\) 网格上的 \eng{HR} 图像，\(S_{t_k}\) 以 contour-refined 位移 \(t_k\) 对 \eng{HR} 网格作亚像素平移重采样，
\(H\) 为 Gaussian \eng{PSF} 卷积（\(\sigma \in [0.2, 0.5]\) \eng{LR} px，见 A.5.2），\(B\) 为 10\,\(\mu\)m 探测器孔径 box 积分，
\(D\) 为 \eng{HR}\(\rightarrow\)\eng{LR} 降采（scale\,=\,2），\(n_k\) 为噪声底 0.0724\,\(^\circ\)C。

仓库中有两套该模型的实现：\eng{EP06} 的矩阵式 forward/adjoint 用于经典重建（\eng{MAP-TV}/\eng{TGV}），
\eng{EP07} 训练 loss 链（blur \(\rightarrow\) block-average \(\rightarrow\) 可选带限）用于深度学习。两者共享 \(H\) 与 \(B \cdot D\) 的核心低通-折叠结构。下文的零空间论证对两者同样成立。

固定 \(t_k\) 后，每帧算子
\[
  A_k = D \cdot B \cdot H \cdot S_{t_k}
\]
是线性的；248 帧 clean 主 session 的 stacked burst 算子写为
\[
  A = \begin{bmatrix} A_1 \\ \vdots \\ A_K \end{bmatrix}, \qquad K=248.
\]
除特别说明外，\(\|\cdot\|\) 为 Euclidean norm。令 \(A = U\Sigma V^\top\) 为完整奇异值分解，右奇异向量为 \(\{v_i\}\)，奇异值为
\[
  \sigma_1 \ge \sigma_2 \ge \cdots \ge \sigma_N \ge 0,\qquad A v_i = \sigma_i u_i.
\]
精确零空间与 \(\epsilon\)-零空间定义为
\[
  \ker(A)=\mathrm{span}\{v_i:\sigma_i=0\},\qquad
  \mathcal N_\epsilon=\mathrm{span}\{v_i:\sigma_i\le\epsilon\}.
\]
观测域损失统一写作 \(\mathcal L(x)=\ell(Ax,y)\)。
若使用 highpass/full-band 等观测侧频带锚，则写作 \(\mathcal L_W(x)=\ell(WAx,Wy)\)，其中 \(W\) 是作用在观测域的有界线性算子。
下列命题只证明固定线性观测算子的性质：它们不证明网络或 \eng{SGD} 一定漂移；它们证明的是，一旦某个漂移分量落在 \(\ker(A)\) 或 \(\mathcal N_\epsilon\) 内，
有界观测侧锚定就只能提供零或趋零的信息与恢复力。

\subsubsection{零空间刻画}

令 \(A = [A_1;\ldots;A_K]\)。stacked 算子的零空间由逐帧零空间交集给出：
\[
  \ker(A)=\bigcap_{k=1}^K \ker(A_k),
\]
且对任意扰动 \(\delta x\)，
\[
  \|A\delta x\|^2 = \sum_{k=1}^K \|A_k\delta x\|^2.
\]
因此，一个单位方向对整个 burst 近似不可观测，当且仅当它在所有帧上的 root-sum-square 响应都很小。
多帧相位多样性只能增加约束：它把单帧条件替换为 \(K\) 个条件的交集，从而可缩小 exact null space，并改善覆盖良好的 alias block 的奇异值；
但它不能恢复被共同的 pre-sampling transfer function \(B H\) 消掉或压到噪声以下的频率。

从机制上看，零空间源于两个来源的叠加。

\textbf{第一，带限衰减。} \(H\) 和 \(B\) 均为低通算子——如 A.1.1 所示，Gaussian \eng{MTF} 在 \(2\times\) Nyquist 处最低仅剩 0.007（\(\sigma = 0.5\)），
叠加 box sinc 后更低。\eng{HR} 网格上超过截止频率的高频分量经 \(A\) 映射后衰减至噪声以下，在数值意义上落入 \(\epsilon\)-零空间。

\textbf{第二，混叠折叠。} 降采算子 \(D\) 将 \eng{HR} 频谱折叠进 \eng{LR} Nyquist 带宽，存在整族高频扰动 \(\delta x\) 使得折叠后各别名分量相消，
严格满足 \(A \cdot \delta x = 0\)。

在周期、shift-invariant 的近似下，这一点可写成 Fourier alias block。对 \(2\times\) 降采，一个 \eng{LR} base frequency \(\omega\) 接收四个 \eng{HR} alias：
\[
  \mathcal A(\omega)=\{\omega+\pi q:q\in\{0,1\}^2\}.
\]
若 \(m(\xi)\) 是 \(B H\) 在 \eng{HR} 频率 \(\xi\) 处的传递系数，则第 \(k\) 帧观测系数为
\[
  \widehat{(A_k\delta x)}(\omega)
  =\sum_{\xi\in\mathcal A(\omega)} m(\xi)e^{-i\xi\cdot t_k}c_\xi.
\]
对应的 stacked alias block 为
\[
  \Phi_\omega(k,\xi)=m(\xi)e^{-i\xi\cdot t_k}.
\]
精确 alias-cancellation 方向就是 \(\ker(\Phi_\omega)\)。更多相位可提高 \(\Phi_\omega\) 的 rank，去掉部分 downsampling-only alias nulls；
但若某个 \(m(\xi)=0\)，该频率是所有帧共有的 transfer zero，任意相位都不可见。若 \(m(\xi)\) 非零但很小，则该方向进入 \(\epsilon\)-零空间，统计上仍不可靠。

直观地说，「细于截止周期的结构改动」与「在所有已采相位下相消的别名组合」对 \(1\times\) 观测不可见或近似不可见。
多帧相位多样性可缩小但不消灭该盲区（25 个相位 bin 的覆盖分析见 A.4.1 和 B.2.3）。

\subsubsection{P1：精确零空间不变性}

\prop{P1}{精确零空间不变性}设 \(\delta x_{\mathrm{null}}\) 满足 \(A \cdot \delta x_{\mathrm{null}} = 0\)，
\(\mathcal{L}(\hat{x}) = \ell(A\hat{x}, y)\) 为任意仅经 \(A\) 接触预测的观测域损失。则对任意 \(\hat{x}\)：
\[
  \mathcal{L}(\hat{x} + \delta x_{\mathrm{null}}) = \ell(A\hat{x} + A\delta x_{\mathrm{null}},\; y) = \ell(A\hat{x},\; y) = \mathcal{L}(\hat{x}).
\]
若 \(\ell\) 对第一自变量可微，则链式法则给出 \(\nabla\mathcal L(\hat{x})=A^\top\nabla_1\ell(A\hat{x},y)\)，因此
\[
  \langle \nabla\mathcal L(\hat{x}),\delta x_{\mathrm{null}}\rangle
  =\langle \nabla_1\ell(A\hat{x},y),A\delta x_{\mathrm{null}}\rangle=0.
\]
此命题的推论是：对 exact null component，任意有限权重的 forward-consistency loss、任意有界 highpass/full-band 频带权重、
以及任意只通过 \(A\) 接触预测的残差范数，都没有可用信息。这正是主文~\S5.2 中 \eng{Stats+HP-FC} / \eng{Stats+Full-FC} 与无锚 \eng{Stats}（run ids: V9B/V9D/v8.1a）漂移曲线重合的机制解释（数字见 D.4.3）。

\subsubsection{P2：\texorpdfstring{\(\epsilon\)}{epsilon}-零空间灵敏度界}

\prop{P2}{\(\epsilon\)-零空间灵敏度}对任意扰动 \(\delta x\)，
\[
  \|A\delta x\|^2=\sum_i \sigma_i^2\langle v_i,\delta x\rangle^2.
\]
若 \(\ell\) 在 \(A\hat{x}\) 与 \(A(\hat{x}+\delta x)\) 之间的线段上对第一自变量是 \(L_\ell\)-Lipschitz，则
\[
  |\mathcal L(\hat{x}+\delta x)-\mathcal L(\hat{x})|
  \le L_\ell\|A\delta x\|.
\]
特别地，若 \(\|\delta x\|=1\) 且 \(\delta x\in\mathcal N_\epsilon\)，则
\[
  \|A\delta x\|\le\epsilon,\qquad
  |\mathcal L(\hat{x}+\delta x)-\mathcal L(\hat{x})|\le L_\ell\epsilon.
\]
对带权观测域损失 \(\mathcal L_W(\hat{x})=\ell(WA\hat{x},Wy)\)，若 \(W\) 有界，则
\[
  |\mathcal L_W(\hat{x}+\delta x)-\mathcal L_W(\hat{x})|
  \le L_\ell\|W\|\,\|A\delta x\|.
\]
因此，\(\|W\|\le1\) 的 normalized highpass/full-band 权重不会比未加权观测损失更敏感；未归一化但有界的 \(W\) 也只能把界放大有限倍。
若 \(A\delta x\) 恰落在 \(W\) 抑制的频带内，灵敏度还会进一步下降，甚至为零。

证明只需把 \(\delta x\) 展开到右奇异向量基中：
\[
  \delta x=\sum_i\langle v_i,\delta x\rangle v_i,\qquad
  A\delta x=\sum_i\sigma_i\langle v_i,\delta x\rangle u_i,
\]
再应用 Lipschitz 不等式与 \(\|WA\delta x\|\le\|W\|\|A\delta x\|\)。
注意，对平方损失而言 Lipschitz 常数是局部的；下一命题直接计算平方损失，不依赖全局 Lipschitz 假设。

\subsubsection{P3：平方损失恢复曲率与有限权重推论}

\prop{P3}{平方观测损失的恢复曲率}对平方数据项
\[
  \mathcal L(\hat{x})=\frac12\|A\hat{x}-y\|^2,
\]
有 \(\nabla\mathcal L(\hat{x})=A^\top(A\hat{x}-y)\)。沿右奇异方向 \(v_i\)，
\[
  \langle\nabla\mathcal L(\hat{x}),v_i\rangle
  =\sigma_i\langle u_i,A\hat{x}-y\rangle
  =\sigma_i^2\langle v_i,\hat{x}\rangle-\sigma_i\langle u_i,y\rangle,
\]
且在直线 \(\hat{x}+\alpha v_i\) 上，
\[
  \frac{d^2}{d\alpha^2}\mathcal L(\hat{x}+\alpha v_i)=\sigma_i^2.
\]
因此，有限权重 \(\lambda\mathcal L\) 在 \(v_i\) 方向上的恢复曲率为 \(\lambda\sigma_i^2\)，随 \(\sigma_i\to0\) 消失。对有界频带加权平方损失
\(\mathcal L_W(\hat{x})=\frac12\|W(A\hat{x}-y)\|^2\)，该方向的曲率为
\[
  \frac{d^2}{d\alpha^2}\mathcal L_W(\hat{x}+\alpha v_i)=\|WAv_i\|^2\le\|W\|^2\sigma_i^2.
\]
所以，有界观测频带权重不能在 \(A\) 的 near-null 方向制造非消失曲率。

\coro{P3}{有限观测权重无法平衡 \(O(1)\) 近零空间力}固定 \(\sigma_i>0\) 的奇异方向 \(v_i\)。
若除有限权重观测项 \(\lambda\mathcal L\) 外，还有一个先验项或网络诱导项 \(G\)，其在 \(v_i\) 方向的局部导数为非零常数 \(g_i\)，
并在所考虑的小段内近似不变，则相对 data-only stationary point 的一维平衡位移必须满足
\[
  |\alpha|=\frac{|g_i|}{\lambda\sigma_i^2}.
\]
若使用带权观测损失，分母变为 \(\lambda\|WAv_i\|^2\)，且 \(\lambda\|WAv_i\|^2\le\lambda\|W\|^2\sigma_i^2\)。
于是，对任意固定有限 \(\lambda\) 与有界 \(W\)，若竞争力 \(g_i=O(1)\)，平衡该力所需的位移会随 \(\sigma_i\to0\) 发散。
反过来说，若人为选择量级为 \(1/\sigma_i^2\) 的无界逆 \eng{MTF} 权重，代数上确实可制造 \(O(1)\) 曲率；
但这不是新的观测信息，而是逆奇异值放大。P4 说明同一放大会同步放大噪声。

\subsubsection{P4：统计不可辨识方差下界}

\prop{P4}{奇异坐标的方差下界}假设
\[
  y=A x^\star+n,\qquad \mathbb E[n]=0,\qquad \mathrm{Cov}(n)=\eta^2 I.
\]
令 \(c_i=\langle v_i,x^\star\rangle\)。对任何线性估计器 \(\hat c_i=a^\top y\)，若它对所有 \(x^\star\) 都无偏估计 \(c_i\)，则
\[
  \mathrm{Var}(\hat c_i)\ge\frac{\eta^2}{\sigma_i^2}.
\]
等号由 \(a=u_i/\sigma_i\) 达到。若 \(\sigma_i=0\)，则不存在对所有 \(x^\star\) 无偏的线性估计器。若噪声为 Gaussian，同一界也是该正交坐标的 Cramér-Rao 下界。

证明要点是，无偏性要求 \(A^\top a=v_i\)。与 \(v_i\) 取内积得
\[
  1=\langle A^\top a,v_i\rangle=\langle a,Av_i\rangle=\sigma_i\langle a,u_i\rangle\le\sigma_i\|a\|,
\]
故 \(\|a\|^2\ge1/\sigma_i^2\)，而 \(\mathrm{Var}(a^\top y)=\eta^2\|a\|^2\)。

P1、P2/P3 与 P4 分别对应三层失效：\(\sigma_i=0\) 时 exact invariance；\(\sigma_i\) 很小时，观测损失变化为 \(O(\sigma_i)\)、平方损失恢复曲率为 \(O(\sigma_i^2)\)；
在带噪声数据中，\(c_i\) 的 Fisher information 也只有 \(\sigma_i^2/\eta^2\)。
因此，near-null 成分必须靠先验、输入证据或选点协议管理，而不能指望另一个有限观测侧锚定项恢复。

\subsubsection{两曲线诊断与直接测量状态}

\textbf{两曲线诊断的充分条件。} 若训练过程中 (i) forward loss 已贴近其噪声地板而无法继续下降，
且 (ii) 真实数据上的 proxy 指标（artifact 与 corr，定义见 A.3）持续单调漂移，
则这强烈指示漂移方向的观测域投影约为零——即漂移主要发生在 \(A\) 的 \(\epsilon\)-零空间内。
该判断是充分条件式的推断，不是恒等式，也不证明网络必然漂移；它只说明一旦实测漂移同时满足这两条曲线特征，有界观测侧锚定缺少可用恢复力。
在 \eng{Stats+HP-FC}（v9b）实验中，forward loss 自 10K 步起平坦于 0.004--0.009，而同期 artifact 从 0.37 单调上爬至 0.65，恰好满足该充分条件。

\textbf{零空间投影的直接测量（待补充 / stretch goal）。} 计划取同一变体相邻 checkpoint 的预测差 \(\delta\hat{x} = \hat{x}_{k+1} - \hat{x}_k\)，
计算 range 分量 \(\delta x_{\mathrm{range}} = A^\dagger A \cdot \delta\hat{x}\)（\(A^\dagger\) 用共轭梯度近似），报告 \(\|\delta x_{\mathrm{range}}\| / \|\delta\hat{x}\|\) 随训练步数的变化。
若该比值在漂移期趋近零，则可将间接诊断升级为直接测量。该项尚未完成，不能作为本文已交付证据。

\subsection{Proxy 反相关的构造性论证}

主文使用 artifact score 与 raw-control correlation 构成的 proxy 对来监测训练漂移。本节论证二者为何存在构造性的反相关关系。

\subsubsection{指标定义}

两个 proxy 共享同一预处理：将 \eng{SR} 输出温度图变换到 highpass 域 \(u = \hat{x} - G_{\sigma_{\mathrm{bg}}}(\hat{x})\)（\(\sigma_{\mathrm{bg}} = 5\) \eng{HR} px），
控制图 \(c\) 为 248 帧原始数据均值的 bicubic 上采样经同一 highpass 变换。

TB-scale 伪影分数 \(\artscore(u)\)（用于训练期 eval 和主文 F3/F4 轨迹图）定义为
\[
  \artscore(u) = \frac{\mathrm{std}(h(u)) + 0.25 \cdot \mathrm{std}(\nabla^2 u)}{\mathrm{std}(u)},\quad h(u)=u-G_{\sigma=1}(u),
\]
即 \(u\) 内部更高频能量的占比。\eng{EP11-harness} 口径使用不同定义（ringing \(+\) 0.25\(\times\)blockiness），数值与 TB-scale 不可比，不得混入同一表格。

Raw-control correlation 定义为 \(\mathrm{corr} = \mathrm{Pearson}(u, c)\)，取全图 finite 像素。

\subsubsection{一阶反相关引理}

两个 proxy 都是同一张 highpass 图 \(u\) 的泛函——artifact 度量 \(u\) 的内部高频能量分布，corr 度量 \(u\) 与固定观测锚 \(c\) 之间的线性一致性。
它们并非独立的证据源，这是反相关论证的前提。但反相关并非无条件成立：只有当扰动比当前残差更高频、且不与固定 raw-control highpass map 对齐时，
artifact 才会一阶上升、corr 才会一阶下降。

为写出条件，令所有向量都已去均值，且 \(c\ne0\)。定义
\[
  L_1=I-G_{\sigma=1},\qquad L_2=\nabla^2,\qquad \beta=0.25,
\]
并令 \(N(u)=\|L_1u\|+\beta\|L_2u\|\)，\(S(u)=\|u\|\)。
于是 artifact proxy 等价于 \(N(u)/S(u)\)（std 的常数因子在比值中抵消），raw-control correlation 为
\[
  \mathrm{corr}(u)=\frac{\langle u,c\rangle}{\|u\|\|c\|}.
\]

\lemmastate{L5}{一阶反相关}令 \(u(\alpha)=u_0+\alpha e\)，并假设 \(\|u_0\|>0\)、\(\|L_1u_0\|>0\)、\(\|L_2u_0\|>0\)，且 \(\rho_0:=\mathrm{corr}(u_0)>0\)。
定义 broadband alignment \(s_0(e)=\dfrac{\langle u_0,e\rangle}{\|u_0\|^2}\)，以及 high-frequency directional growth
\[
  h_0(e)=
  \frac{
  \frac{\langle L_1u_0,L_1e\rangle}{\|L_1u_0\|}
  +\beta\frac{\langle L_2u_0,L_2e\rangle}{\|L_2u_0\|}
  }{
  \|L_1u_0\|+\beta\|L_2u_0\|
  }.
\]
若 \(h_0(e)>s_0(e)\)，则
\[
  \left.\frac{d}{d\alpha}\artscore(u_0+\alpha e)\right|_{\alpha=0}>0.
\]
若进一步满足 \(\dfrac{\langle e,c\rangle}{\|u_0\|\|c\|}<\rho_0\,s_0(e)\)，则
\[
  \left.\frac{d}{d\alpha}\mathrm{corr}(u_0+\alpha e)\right|_{\alpha=0}<0.
\]
也就是说，对边缘增强式的 style perturbation，只要它相对当前 \(u_0\) 增加更高频能量，且与固定 raw-control map \(c\) 的对齐很弱，
artifact 会一阶上升而 corr 会一阶下降。

证明来自范数一阶导数。对任意非零可微路径 \(q(\alpha)\)，\(\left.\frac{d}{d\alpha}\|q(\alpha)\|\right|_{\alpha=0}=\frac{\langle q(0),q'(0)\rangle}{\|q(0)\|}\)。
代入 \(L_1u(\alpha)\)、\(L_2u(\alpha)\) 与 \(u(\alpha)\)，得到
\[
  N'(0)=
  \frac{\langle L_1u_0,L_1e\rangle}{\|L_1u_0\|}
  +\beta\frac{\langle L_2u_0,L_2e\rangle}{\|L_2u_0\|},
  \qquad
  S'(0)=\frac{\langle u_0,e\rangle}{\|u_0\|}.
\]
由 \(\artscore(u)=N(u)/S(u)\) 可知，\(\artscore\) 的一阶导数为正当且仅当 \(N'(0)/N(u_0)>S'(0)/S(u_0)\)，即 \(h_0(e)>s_0(e)\)。对 correlation，
\[
  \left.\frac{d}{d\alpha}\mathrm{corr}(u_0+\alpha e)\right|_{\alpha=0}
  =
  \frac{\langle e,c\rangle}{\|u_0\|\|c\|}
  -\rho_0\frac{\langle u_0,e\rangle}{\|u_0\|^2},
\]
于是上式小于零正是第二个条件。

若 \(e\perp u_0\) 且 \(e\perp c\)，corr 的一阶导数为零而非负。此时若 \(\rho_0>0\)，下降只在二阶出现：
\[
  \mathrm{corr}(u_0+\alpha e)
  =\rho_0\left(1+\alpha^2\frac{\|e\|^2}{\|u_0\|^2}\right)^{-1/2}.
\]
因此，主文中的「构造性反相关」应理解为满足上述方向条件的局部诊断，而不是任意扰动方向上的恒等关系。

由此得到三条推论。第一，proxy 对是漂移温度计与选点准则，不是两个可联合优化的分数。
第二，跨 input-mode 的横向比较无效：hybrid drizzle 输入合法携带更多高频能量，其 artifact 基线天然更高。
第三，应关注训练轨迹的走向而非端点数值。

\subsubsection{假设与适用边界}

\begin{itemize}
  \item \textbf{固定线性算子。} 上述命题假设 \(t_k\) 已固定，因此 \(S_{t_k}\) 与 \(A_k=D B H S_{t_k}\) 是线性算子。若把位移与 \(x\) 联合优化，问题不再是本文证明的固定线性观测算子。
  \item \textbf{有限维离散化。} \eng{SVD} 命题对实现中的有限矩阵精确成立；Fourier/MTF 解释还额外假设周期或内部区域近似 shift-invariant。非周期边界会扰动 Fourier 对角化，但不影响代数 \eng{SVD} 结论。
  \item \textbf{相位多样性。} stacked exact null space 是 \(\bigcap_k\ker(A_k)\)。更多相位可去掉部分单帧混叠零空间，却不能去掉 \(B H\) 的共同 transfer zeros，也不能让极小 \eng{MTF} 方向变得统计可靠。
  \item \textbf{有界观测侧带权。} 「频带锚定不能补救」指的是有界 \(W\) 与固定有限损失权重 \(\lambda\)。无界逆滤波或按 \(1/\sigma_i^2\) 放大的权重可在代数上造曲率，但会按同一逆奇异值放大噪声，不属于有限观测锚定主张。
  \item \textbf{损失正则性。} P2 使用 \(\ell\) 在相关残差范围内的 Lipschitz 界；平方损失不是全局 Lipschitz，因此 P2 对平方损失只作局部解释，P3 给出精确平方损失计算。
  \item \textbf{噪声模型。} P4 的线性无偏估计界假设零均值、协方差为 \(\eta^2 I\) 的噪声。若噪声相关，应改用 whitened operator 或 Fisher information \(A^\top C_n^{-1}A\)。
  \item \textbf{不声明神经网络动力学。} 这些命题不证明 \eng{UNet} 或 \eng{SGD} 必然漂移，也不分析非凸训练动力学；它们只说明，一旦漂移分量位于 exact/near null space，有界观测域损失在该分量上提供零或趋零的信息与恢复力。
  \item \textbf{proxy 反相关是条件性的。} L5 要求扰动相对 \(u_0\) 更高频、边缘增强或至少增加分母能量，并且近似正交于固定 raw-control highpass map \(c\)。若条件不满足，一阶符号可为零或改变。
\end{itemize}

\subsection{FRC 方法学}

本节详述 Fourier Ring Correlation（\eng{FRC}）的实验方案与全部控制组结果。

\subsubsection{分层分半构造}

\eng{FRC} 的核心思路是将 248 帧 clean set 分为两个子集 A、B，各自独立重建后计算频域相关。
为保证分半不引入系统性偏差，采用相位分层（phase-stratified）设计：将每帧的 stage command 经坐标变换映射为亚像素相位 \(\phi = \mathrm{mod}(\text{shift}, 1)\)，
按 \(5 \times 5 = 25\) 个 bin 分层（每 bin 7--13 帧），在每个 bin 内随机置换后奇偶交替分配给 A/B（约束 \(|N_A - N_B| \le 1\)）。
A 和 B 各自以 bilinear drizzle 重建到 \(5\times\) 诊断网格（hr\_pitch \(=\) 2\,\(\mu\)m），空 bin 填全局均值。
使用 seeds \{42, 123, 456\} 重复三次，主曲线取逐 ring nanmean。窗口函数为 Tukey \(\alpha = 0.25\)，边缘裁剪 16 \eng{LR} px。

\eng{FRC} 按环状积分计算：
\[
  \mathrm{FRC}(r) = \frac{\sum_{|\mathbf{f}| \in r} \mathrm{Re}\bigl(F_A(\mathbf{f}) \cdot \overline{F_B(\mathbf{f})}\bigr)}{\sqrt{\sum_{|\mathbf{f}| \in r} |F_A(\mathbf{f})|^2 \cdot \sum_{|\mathbf{f}| \in r} |F_B(\mathbf{f})|^2}}.
\]
cutoff 判据采用 \(1/7\) 常数阈值（0.142857）和 half-bit 曲线，取第一个 \eng{FRC} 低于阈值的 ring 对应周期。

\subsubsection{主结果与控制组}

\(1/7\) cutoff 为 17.03\,\(\mu\)m（3-seed 均值曲线），逐 seed 分别为 16.17/16.17/17.03\,\(\mu\)m（std \(=\) 0.50\,\(\mu\)m）；half-bit 判据给出相同值。
本文仅主张 17.0\,\(\mu\)m：超过 20\,\(\mu\)m 分辨率的相干信息确实存在，但低于 11--14\,\(\mu\)m 的理论期望。

四个控制组的设计与结果如下。正控制（单帧 bicubic 加噪）的 cutoff 为 13.58\,\(\mu\)m，未达到「明显差于 main」的预期，
原因是单帧插值的平滑频谱在低噪声下自相关偏高；负控制（shift-shuffle）在 8--12\,\(\mu\)m 段中位 \eng{FRC} 仍有 0.504，
因为置换后仍共享场景低频与网格结构，部分失效。漂移控制（前/后半按采集序分别重建）cutoff 退化至 26.20\,\(\mu\)m，
符合预期——时间间隔放大了热漂移的影响。zero-coverage 统计显示空 bin 均值 27.2\%、最大 36.2\%，表明 \(5\times\) 网格存在欠覆盖。

\subsubsection{10--12 \texorpdfstring{\(\mu\)}{u}m 反弹}

在 10--12\,\(\mu\)m 周期段，\eng{FRC} 出现异常高值，需要谨慎判读。关键频带的对比数据如下。

\begin{table}[ht]
  \centering
  \caption{关键周期段的 \eng{FRC}：main 与三个控制组对比}
  \label{tab:frc-band}
  \small
  \begin{tabular}{rcccc}
    \toprule
    周期 (\(\mu\)m) & main & bicubic 正控 & shuffle 负控 & drift 控 \\
    \midrule
    16 & 0.138 & 0.691 & \(-0.012\) & 0.117 \\
    12 & 0.593 & 0.002 & \(-0.095\) & 0.578 \\
    10 & 0.935 & 0.012 & 0.906 & 0.887 \\
    8  & 0.545 & \(-0.023\) & 0.390 & 0.345 \\
    \bottomrule
  \end{tabular}
\end{table}

10--12\,\(\mu\)m 的高 \eng{FRC} 同时出现在负控制（10\,\(\mu\)m 处 0.906）和漂移控制（0.887）中，
说明该反弹由 coverage/lattice 结构与漂移驱动，而非真实分辨率信息。因此不作分辨率证据使用。

\subsubsection{MAP-TV 前后对照}

以 \eng{MAP-TV}（\(\sigma = 0.2\), \(\lambda = 10^{-3}\)）对 drizzle 输出做去卷积后，split-half \eng{FRC} 在 20--10\,\(\mu\)m 段从 0.319/0.088/0.053/0.575/0.893 提升至 0.976/0.965/0.955/0.947/0.934。
该提升反映的是分半一致性（稳定性）的改善，而非光学分辨率本身的提升。

\subsection{标定不确定度传播}

\subsubsection{旋转角 \texorpdfstring{\(\theta\)}{theta} 的误差传播}

坐标到位移的映射为
\[
  d_x = \frac{X\cos\theta + Y\sin\theta}{10}, \quad d_y = \frac{-X\sin\theta + Y\cos\theta}{10}
\]
（单位 \(\mu\)m\(\rightarrow\)px）。对 \(\theta\) 求导，以最大步进 \(X = 40\)\,\(\mu\)m、\(\delta\theta = 0.1° = 1.745\)\,mrad 为例：
\(\delta d_x \approx 0.0052\)\,px，\(\delta d_y \approx 0.0047\)\,px，合成约 0.007\,px。
位移幅值 \(|s| = X/10 = 4.0\)\,px 与 \(\theta\) 一阶无关——\(\theta\) 误差表现为方向偏差而非步长误差。
该传播量比 alignment 残差（Chamfer 0.134\,px，见 B.2.2）低一个量级，不是误差预算的主项。

\eng{AVI} 的独立方向验证（gradient-NCC 合并估计 \(\theta = 47.14°\)，95\% CI [46.36°, 47.92°]）覆盖了标定值 47.6°，
但 X/Y 子组存在约 3° 系统差异（X 48.70° vs Y 45.63°），且 \eng{AVI} 为 8-bit 渲染视频，仅作一致性检验，不替换配置。

\subsubsection{PSF \texorpdfstring{\(\sigma\)}{sigma} 区间}

三条独立标定路线给出了不一致的结果：Route A（forward 残差）\(\sigma = 0.226\)\,px [0.208, 0.240]；
Route B（ESF 拟合）\(\sigma = 1.129\)\,px [1.041, 1.215]；Route C（joint hold-out）\(\sigma = 0.119\)\,px。
三者 spread 达 1.01\,px，远超 0.05\,px 容差。

\eng{EP15 M3} 仲裁澄清了分歧的物理根源：Route B 的 1.129\,px 实际上是 \eng{PSF} 与热/几何边缘宽度的卷积表观值，满足 \(\sigma_{\mathrm{total}}^2 = \sigma_{\mathrm{PSF}}^2 + w_{\mathrm{edge}}^2\)。
这意味着边缘表观宽度不能直接等同于 \eng{PSF}，同时也为反卷积的激进度设定了物理上限——即使完全去除光学 \eng{PSF}，热扩散造成的边缘宽度仍然存在。
综合 \eng{FRC} 形状拟合（\(\sigma = 0.2\) 最优）与多边缘 ESF 上界（0.546\,px），最终采纳区间为 \(\sigma \in [0.2, 0.5]\) \eng{LR} px。

该区间的传播后果显著：\(2\times\) Nyquist 处 \eng{MTF} 在区间内变化 65 倍（0.454\(\rightarrow\)0.007），因此所有依赖 \eng{PSF} 的计算均须把 \(\sigma\) 作为扫描参数。

\subsubsection{误差预算汇总}

\begin{table}[ht]
  \centering
  \caption{误差预算汇总}
  \label{tab:error-budget}
  \small
  \begin{tabular}{p{0.24\linewidth}p{0.32\linewidth}p{0.34\linewidth}}
    \toprule
    误差源 & 量级 & 对结论的影响 \\
    \midrule
    \(\theta\) \(\pm\)0.1° & \(\sim\)0.007\,px @40\,\(\mu\)m & 可忽略（不足 alignment 残差的 5\%） \\
    alignment 残差 & Chamfer 0.134\,px（refined） & 主要几何误差项；E3 消融显示改善端到端 corr \(+\)0.11（见 D.5.3） \\
    PSF \(\sigma\) 区间 & [0.2, 0.5]\,px \(\rightarrow\) \(2\times\) MTF 65\(\times\) 跨度 & 反卷积与合成均按区间扫描 \\
    噪声底 & 0.0724\,\(^\circ\)C & 可恢复对比度的下限 \\
    热漂移 & session 内 \(-0.60\)\,\(^\circ\)C；跨 session 中位 2.91\,\(^\circ\)C & session 门控 \(+\) 分层 split 的设计动机（见 B.1、B.3） \\
    \bottomrule
  \end{tabular}
\end{table}
